\documentclass[11pt,oneside]{article}

\pdfminorversion=7
\usepackage{QuizStyle}

\hypersetup{
    pdftitle={20261006 Quiz for MATH230 Probability \& Statistics},
    pdfauthor={Jungwoo Kim},
    pdfsubject={Quiz for MATH230 Probability \& Statistics},
    pdfcreator={LaTeX}
}

\begin{document}

\quiztitle{20261006 Quiz}

\begin{exercise}{4.34}
Let $X$ be a random variable with the following probability distribution:
\[
\begin{array}{c@{\qquad}ccc}
\toprule
x & -2 & 3 & 5 \\
\midrule
f(x) & 0.3 & 0.3 & 0.4 \\
\bottomrule
\end{array}
\]
Find the standard deviation of $X$.
\end{exercise}

\begin{solution}
By Theorem 4.2,
\[
\begin{aligned}
    \Exp[X]
    &=(-2)(0.3)+(3)(0.3)+(5)(0.4)=\frac{23}{10},\\
    \Exp[X^2]
    &=(-2)^2(0.3)+3^2(0.3)+5^2(0.4)=\frac{139}{10}.
\end{aligned}
\]
Hence,
\[
    \Var(X)=\frac{139}{10}-\left(\frac{23}{10}\right)^2
    =\frac{861}{100},
\]
so
\[
    \boxed{\sigma_X=\frac{\sqrt{861}}{10}\approx2.93}.
\]
\end{solution}

\begin{exercise}{4.35}
The random variable $X$, representing the number of errors per 100 lines of software code, has the following probability distribution:
\[
\begin{array}{c@{\qquad}ccccc}
\toprule
x & 2 & 3 & 4 & 5 & 6 \\
\midrule
f(x) & 0.01 & 0.25 & 0.4 & 0.3 & 0.04 \\
\bottomrule
\end{array}
\]
Using Theorem 4.2 on page 141, find the variance of $X$.
\end{exercise}

\begin{solution}
Theorem 4.2 gives
\[
    \Var(X)=\Exp[X^2]-\bigl(\Exp[X]\bigr)^2.
\]
Here,
\[
\begin{aligned}
    \Exp[X]
    &=(2)(0.01)+(3)(0.25)+(4)(0.4)+(5)(0.3)+(6)(0.04)
      =4.11,\\
    \Exp[X^2]
    &=2^2(0.01)+3^2(0.25)+4^2(0.4)+5^2(0.3)+6^2(0.04)
      =17.63.
\end{aligned}
\]
Therefore,
\[
    \boxed{\Var(X)=17.63-(4.11)^2=0.7379\approx0.74}.
\]
\end{solution}

\begin{exercise}{4.36}
Suppose that the probabilities are $0.4$, $0.3$, $0.2$, and $0.1$, respectively, that $0$, $1$, $2$, or $3$ power failures will strike a certain subdivision in any given year. Find the mean and variance of the random variable $X$ representing the number of power failures striking this subdivision.
\end{exercise}

\begin{solution}
The distribution is
\[
    p_X(0)=0.4,\quad p_X(1)=0.3,\quad p_X(2)=0.2,\quad p_X(3)=0.1.
\]
Thus,
\[
\begin{aligned}
    \Exp[X]&=(1)(0.3)+(2)(0.2)+(3)(0.1)=1,\\
    \Exp[X^2]&=(1)^2(0.3)+(2)^2(0.2)+(3)^2(0.1)=2.
\end{aligned}
\]
By Theorem 4.2,
\[
    \Var(X)=2-1^2=1.
\]
Hence,
\[
    \boxed{\mu_X=1,\qquad \Var(X)=1}.
\]
\end{solution}

\begin{exercise}{4.37}
A dealer's profit, in units of $\$5000$, on a new automobile is a random variable $X$ having the density function given in Exercise 4.12 on page 137. Find the variance of $X$.
\end{exercise}

\begin{solution}
From Exercise 4.12,
\[
    f(x)=
    \begin{cases}
        2(1-x), & 0<x<1,\\
        0, & \text{elsewhere}.
    \end{cases}
\]
Direct integration gives
\[
\begin{aligned}
    \Exp[X]
    &=\int_0^1 2x(1-x)\,\dd x=\frac13,\\
    \Exp[X^2]
    &=\int_0^1 2x^2(1-x)\,\dd x=\frac16.
\end{aligned}
\]
Therefore,
\[
    \boxed{\Var(X)=\frac16-\left(\frac13\right)^2=\frac1{18}}.
\]
Since $X$ is measured in units of $\$5000$, the variance of the actual profit in dollars is
\[
    \frac1{18}(5000)^2\ \text{dollars}^2.
\]
\end{solution}

\begin{exercise}{4.40}
Referring to Exercise 4.14 on page 137, find $\Var(g(X))$ for the function $g(X)=3X+4$.
\end{exercise}

\begin{solution}
From Exercise 4.14,
\[
    f(x)=
    \begin{cases}
        \dfrac{2(x+2)}{5}, & 0<x<1,\\
        0, & \text{elsewhere}.
    \end{cases}
\]
By Theorem 4.3,
\[
    \Var(g(X))=\Exp[g(X)^2]-\bigl(\Exp[g(X)]\bigr)^2.
\]
The required moments are
\[
\begin{aligned}
    \Exp[g(X)]
    &=\int_0^1 (3x+4)\frac{2(x+2)}5\,\dd x
      =\frac{28}{5},\\
    \Exp[g(X)^2]
    &=\int_0^1 (3x+4)^2\frac{2(x+2)}5\,\dd x
      =\frac{321}{10}.
\end{aligned}
\]
Therefore,
\[
    \boxed{\Var(g(X))
    =\frac{321}{10}-\left(\frac{28}{5}\right)^2
    =\frac{37}{50}=0.74}.
\]
\end{solution}

\begin{exercise}{4.41}
Find the standard deviation of the random variable $h(X)=(3X+1)^2$ in Exercise 4.17 on page 138.
\end{exercise}

\begin{solution}
From Exercise 4.17,
\[
\begin{array}{c@{\qquad}ccc}
\toprule
x & -3 & 6 & 9 \\
\midrule
p_X(x) & \frac16 & \frac12 & \frac13 \\
\bottomrule
\end{array}
\]
and
\[
    h(-3)=64,\qquad h(6)=361,\qquad h(9)=784.
\]
By Theorem 4.3,
\[
\begin{aligned}
    \Exp[h(X)]
    &=64\left(\frac16\right)
      +361\left(\frac12\right)
      +784\left(\frac13\right)
      =\frac{905}{2},\\
    \Exp[h(X)^2]
    &=64^2\left(\frac16\right)
      +361^2\left(\frac12\right)
      +784^2\left(\frac13\right)
      =\frac{541457}{2}.
\end{aligned}
\]
Thus,
\[
    \Var(h(X))
    =\frac{541457}{2}-\left(\frac{905}{2}\right)^2
    =\frac{263889}{4},
\]
and therefore
\[
    \boxed{\sigma_{h(X)}=\frac{\sqrt{263889}}{2}\approx256.85}.
\]
\end{solution}

\begin{exercise}{4.43}
The length of time, in minutes, for an airplane to obtain clearance for takeoff at a certain airport is a random variable $Y=3X-2$, where $X$ has the density function
\[
    f(x)=
    \begin{cases}
        \dfrac14 e^{-x/4}, & x>0,\\
        0, & \text{elsewhere}.
    \end{cases}
\]
Find the mean and variance of the random variable $Y$.
\end{exercise}

\begin{solution}
Since $Y=g(X)=3X-2$, we use the expectation formulas for a function of $X$ directly:
\[
\begin{aligned}
    \Exp[Y]
    &=\int_0^\infty (3x-2)\frac14e^{-x/4}\,\dd x
      =10,\\
    \Exp[Y^2]
    &=\int_0^\infty (3x-2)^2\frac14e^{-x/4}\,\dd x
      =244.
\end{aligned}
\]
\Needspace{6\baselineskip}
Therefore,
\[
    \Var(Y)=244-10^2=144.
\]
Hence,
\[
    \boxed{\mu_Y=10,\qquad \Var(Y)=144}.
\]
\end{solution}

\Needspace{19\baselineskip}
\begin{exercise}{4.44}
Find the covariance of the random variables $X$ and $Y$ of Exercise 3.39 on page 125.
\end{exercise}

\begin{solution}
In Exercise 3.39, four pieces are selected from $3$ oranges, $2$ apples, and $3$ bananas, with $X$ and $Y$ denoting the numbers of oranges and apples selected. Thus
\[
    p_{X,Y}(x,y)
    =\frac{\binom{3}{x}\binom{2}{y}\binom{3}{4-x-y}}{\binom{8}{4}}
\]
for the admissible values of $x$ and $y$. The marginal distributions give
\[
\begin{aligned}
    p_X(0)&=\frac1{14}, & p_X(1)&=\frac37,
    & p_X(2)&=\frac37, & p_X(3)&=\frac1{14},\\
    p_Y(0)&=\frac3{14}, & p_Y(1)&=\frac47,
    & p_Y(2)&=\frac3{14}.
\end{aligned}
\]
Hence,
\[
    \Exp[X]=\frac32,
    \qquad
    \Exp[Y]=1.
\]
Also,
\[
\begin{aligned}
    \Exp[XY]
    &=\sum_x\sum_y xy\,p_{X,Y}(x,y)\\
    &=\frac{18+18+36+12+6}{70}
      =\frac97.
\end{aligned}
\]
Therefore, by Theorem 4.4,
\[
    \boxed{\Cov(X,Y)
    =\Exp[XY]-\Exp[X]\Exp[Y]
    =\frac97-\frac32
    =-\frac3{14}}.
\]
\end{solution}

\begin{exercise}{4.45}
Find the covariance of the random variables $X$ and $Y$ of Exercise 3.49 on page 126.
\end{exercise}

\begin{solution}
From Exercise 3.49, the joint distribution is
\[
\begin{array}{c@{\qquad}ccc}
\toprule
p_{X,Y}(x,y) & \multicolumn{3}{c}{x} \\
\cmidrule(lr){2-4}
y & 1 & 2 & 3 \\
\midrule
1 & 0.05 & 0.05 & 0.10 \\
3 & 0.05 & 0.10 & 0.35 \\
5 & 0.00 & 0.20 & 0.10 \\
\bottomrule
\end{array}
\]
so the marginal distributions are
\[
\begin{aligned}
    p_X(1)&=0.10, & p_X(2)&=0.35, & p_X(3)&=0.55,\\
    p_Y(1)&=0.20, & p_Y(3)&=0.50, & p_Y(5)&=0.30.
\end{aligned}
\]
Thus,
\[
    \Exp[X]=\frac{49}{20},
    \qquad
    \Exp[Y]=\frac{16}{5},
\]
\Needspace{10\baselineskip}
The mixed moment is
\[
\begin{aligned}
    \Exp[XY]
    &=\sum_x\sum_y xy\,p_{X,Y}(x,y)\\
    &=(1)(1)(0.05)+(2)(1)(0.05)+(3)(1)(0.10)\\
    &\qquad +(1)(3)(0.05)+(2)(3)(0.10)+(3)(3)(0.35)\\
    &\qquad +(2)(5)(0.20)+(3)(5)(0.10)\\
    &=\frac{157}{20}.
\end{aligned}
\]
Therefore, by Theorem 4.4,
\[
    \boxed{\Cov(X,Y)
    =\frac{157}{20}-\left(\frac{49}{20}\right)\left(\frac{16}{5}\right)
    =\frac1{100}=0.01}.
\]
\end{solution}

\begin{exercise}{4.48}
Given a random variable $X$, with standard deviation $\sigma_X$, and a random variable $Y=a+bX$, show that if $b<0$, the correlation coefficient $\rho_{X,Y}=-1$, and if $b>0$, $\rho_{X,Y}=1$.
\end{exercise}

\begin{solution}
Assume $\sigma_X>0$, as required for the correlation coefficient to be defined. By Theorem 4.5,
\[
    \Exp[Y]=a+b\Exp[X].
\]
We have
\[
    Y-\Exp[Y]=b\bigl(X-\Exp[X]\bigr).
\]
Hence, directly from the definitions,
\[
\begin{aligned}
    \Cov(X,Y)
    &=b\Var(X)=b\sigma_X^2,\\
    \sigma_Y
    &=|b|\sigma_X.
\end{aligned}
\]
Therefore,
\[
    \rho_{X,Y}
    =\frac{\Cov(X,Y)}{\sigma_X\sigma_Y}
    =\frac{b}{|b|}
    =
    \begin{cases}
        -1, & b<0,\\
        1, & b>0.
    \end{cases}
\]
\end{solution}

\begin{exercise}{4.52}
Random variables $X$ and $Y$ follow a joint distribution
\[
    f(x,y)=
    \begin{cases}
        2, & 0<x\leq y<1,\\
        0, & \text{otherwise}.
    \end{cases}
\]
Determine the correlation coefficient between $X$ and $Y$.
\end{exercise}

\begin{solution}
Over the support $0<x\leq y<1$,
\[
\begin{aligned}
    \Exp[X]
    &=\int_0^1\int_0^y 2x\,\dd x\,\dd y=\frac13,
    &
    \Exp[Y]
    &=\int_0^1\int_0^y 2y\,\dd x\,\dd y=\frac23,\\
    \Exp[X^2]
    &=\int_0^1\int_0^y 2x^2\,\dd x\,\dd y=\frac16,
    &
    \Exp[Y^2]
    &=\int_0^1\int_0^y 2y^2\,\dd x\,\dd y=\frac12,\\
    \Exp[XY]
    &=\int_0^1\int_0^y 2xy\,\dd x\,\dd y=\frac14.
\end{aligned}
\]
Thus,
\[
    \Var(X)=\Var(Y)=\frac1{18},
    \qquad
    \Cov(X,Y)=\frac14-\left(\frac13\right)\left(\frac23\right)=\frac1{36}.
\]
Therefore,
\[
    \boxed{\rho_{X,Y}
    =\frac{1/36}{\sqrt{(1/18)(1/18)}}
    =\frac12}.
\]
\end{solution}

\begin{exercise}{4.57}
Let $X$ be a random variable with the following probability distribution:
\[
\begin{array}{c@{\qquad}ccc}
\toprule
x & -3 & 6 & 9 \\
\midrule
f(x) & \frac16 & \frac13 & \frac12 \\
\bottomrule
\end{array}
\]
Find $\Exp[X]$ and $\Exp[X^2]$ and then, using these values, evaluate $\Exp[(2X+1)^2]$.
\end{exercise}

\begin{solution}
Directly from the distribution,
\[
\begin{aligned}
    \Exp[X]
    &=(-3)\left(\frac16\right)
      +6\left(\frac13\right)
      +9\left(\frac12\right)
      =6,\\
    \Exp[X^2]
    &=9\left(\frac16\right)
      +36\left(\frac13\right)
      +81\left(\frac12\right)
      =54.
\end{aligned}
\]
Thus,
\[
    \boxed{\Exp[X]=6,\qquad \Exp[X^2]=54}.
\]
Since $(2X+1)^2=4X^2+4X+1$, Theorem 4.6 gives
\[
    \boxed{\Exp[(2X+1)^2]
    =4\Exp[X^2]+4\Exp[X]+1
    =241}.
\]
\end{solution}

\Needspace{22\baselineskip}
\begin{exercise}{4.58}
The total time, measured in units of 100 hours, that a teenager runs their hair dryer over a period of one year is a continuous random variable $X$ that has the density function
\[
    f(x)=
    \begin{cases}
        x, & 0<x<1,\\
        2-x, & 1\leq x<2,\\
        0, & \text{elsewhere}.
    \end{cases}
\]
Use Theorem 4.6 to evaluate the mean of the random variable $Y=60X^2+48X$, where $Y$ is equal to the number of kilowatt hours expended annually.
\end{exercise}

\begin{solution}
First,
\[
\begin{aligned}
    \Exp[X]
    &=\int_0^1 x^2\,\dd x+\int_1^2 x(2-x)\,\dd x=1,\\
    \Exp[X^2]
    &=\int_0^1 x^3\,\dd x+\int_1^2 x^2(2-x)\,\dd x=\frac76.
\end{aligned}
\]
By Theorem 4.6,
\[
    \Exp[Y]
    =60\Exp[X^2]+48\Exp[X]
    =60\left(\frac76\right)+48
    =\boxed{118}.
\]
Thus, the mean annual energy expenditure is $118$ kilowatt hours.
\end{solution}

\begin{exercise}{4.60}
Suppose that $X$ and $Y$ are independent random variables having the joint probability distribution
\[
\begin{array}{c@{\qquad}cc}
\toprule
f(x,y) & \multicolumn{2}{c}{x} \\
\cmidrule(lr){2-3}
y & 2 & 4 \\
\midrule
1 & 0.15 & 0.10 \\
3 & 0.25 & 0.25 \\
5 & 0.15 & 0.10 \\
\bottomrule
\end{array}
\]
Find
\begin{exercisesubparts}
    \item $\Exp[2X-3Y]$;
    \item $\Exp[XY]$.
\end{exercisesubparts}
\end{exercise}

\begin{solution}
The marginal means are
\[
    \Exp[X]=(2)(0.55)+(4)(0.45)=\frac{29}{10},
    \qquad
    \Exp[Y]=(1)(0.25)+(3)(0.50)+(5)(0.25)=3.
\]

\noindent\textup{(a)} By Theorems 4.5 and 4.7,
\[
    \boxed{\Exp[2X-3Y]
    =2\Exp[X]-3\Exp[Y]
    =-\frac{16}{5}}.
\]

\medskip
\noindent\textup{(b)} Using the displayed joint distribution directly,
\[
\begin{aligned}
    \Exp[XY]
    &=\sum_x\sum_y xyf(x,y)\\
    &=(2)(1)(0.15)+(4)(1)(0.10)
      +(2)(3)(0.25)+(4)(3)(0.25)\\
    &\qquad +(2)(5)(0.15)+(4)(5)(0.10)\\
    &=\boxed{\frac{87}{10}}.
\end{aligned}
\]
\end{solution}

\Needspace{10\baselineskip}
\begin{exercise}{4.61}
Use Theorem 4.7 to evaluate $\Exp[2XY^2-X^2Y]$ for the joint probability distribution shown in Table 3.1 on page 116.
\end{exercise}

\begin{solution}
By Theorem 4.7,
\[
    \Exp[2XY^2-X^2Y]
    =2\Exp[XY^2]-\Exp[X^2Y].
\]
Both terms contain a factor $XY$, so every cell with $x=0$ or $y=0$ contributes $0$. From Table 3.1, the only remaining nonzero cell is
\[
    p_{X,Y}(1,1)=\frac3{14}.
\]
Hence,
\[
    \Exp[XY^2]=\frac3{14},
    \qquad
    \Exp[X^2Y]=\frac3{14},
\]
and therefore
\[
    \boxed{\Exp[2XY^2-X^2Y]=\frac3{14}}.
\]
\end{solution}

\begin{exercise}{4.65}
Let $X$ represent the number that occurs when a red die is tossed and $Y$ the number that occurs when a green die is tossed. Find
\begin{exercisesubparts}
    \item $\Exp[2X+Y]$;
    \item $\Exp[X-2Y]$;
    \item $\Exp[2XY]$.
\end{exercisesubparts}
\end{exercise}

\begin{solution}
For each die,
\[
    \Exp[X]=\Exp[Y]=\frac16\sum_{k=1}^6 k=\frac72.
\]
Therefore, by Theorems 4.5 and 4.7,
\[
\begin{aligned}
    \textup{(a)}\quad
    \Exp[2X+Y]
    &=2\Exp[X]+\Exp[Y]
      =\boxed{\frac{21}{2}},\\[0.4em]
    \textup{(b)}\quad
    \Exp[X-2Y]
    &=\Exp[X]-2\Exp[Y]
      =\boxed{-\frac72}.
\end{aligned}
\]
For part (c), the two die outcomes are independent, so their joint pmf is $1/36$ on $\{1,\ldots,6\}^2$. Thus
\[
\begin{aligned}
    \Exp[XY]
    &=\frac1{36}
      \left(\sum_{x=1}^6x\right)
      \left(\sum_{y=1}^6y\right)
      =\frac{49}{4},\\
    \textup{(c)}\quad
    \Exp[2XY]
    &=\boxed{\frac{49}{2}}.
\end{aligned}
\]
\end{solution}

\begin{exercise}{4.69}
Consider Review Exercise 3.77 on page 128. The random variables $X$ and $Y$ represent the number of vehicles that arrive at two separate street corners during a certain 2-minute period in the day. The joint distribution is
\[
    f(x,y)=\left(\frac{1}{4^{x+y}}\right)\left(\frac{9}{16}\right),
\]
for $x=0,1,2,\ldots$ and $y=0,1,2,\ldots$.
\begin{exercisesubparts}
    \item Give $\Exp[X]$, $\Exp[Y]$, $\Var(X)$, and $\Var(Y)$.
    \item Consider $Z=X+Y$, the sum of the two. Find $\Exp[Z]$ and $\Var(Z)$.
\end{exercisesubparts}
\end{exercise}

\begin{solution}
The joint pmf factors as
\[
    f(x,y)
    =\left(\frac34\left(\frac14\right)^x\right)
     \left(\frac34\left(\frac14\right)^y\right),
\]
so $X$ and $Y$ are independent and have the same marginal pmf.

\noindent\textup{(a)} For $|r|<1$,
\[
    \sum_{k=0}^{\infty}kr^k=\frac{r}{(1-r)^2},
    \qquad
    \sum_{k=0}^{\infty}k^2r^k=\frac{r(1+r)}{(1-r)^3}.
\]
With $r=1/4$,
\[
\begin{aligned}
    \Exp[X]=\Exp[Y]
    &=\frac34\sum_{k=0}^{\infty}k\left(\frac14\right)^k
      =\frac13,\\
    \Exp[X^2]=\Exp[Y^2]
    &=\frac34\sum_{k=0}^{\infty}k^2\left(\frac14\right)^k
      =\frac59.
\end{aligned}
\]
Therefore,
\[
    \boxed{\Exp[X]=\Exp[Y]=\frac13,
    \qquad
    \Var(X)=\Var(Y)=\frac49}.
\]

\medskip
\noindent\textup{(b)} By Theorem 4.7,
\[
    \Exp[Z]=\Exp[X+Y]=\Exp[X]+\Exp[Y]=\frac23.
\]
For the mixed moment, use the factorized joint pmf directly:
\[
\begin{aligned}
    \Exp[XY]
    &=\sum_x\sum_y xy
      \left(\frac34\left(\frac14\right)^x\right)
      \left(\frac34\left(\frac14\right)^y\right)\\
    &=\left(\frac34\sum_{x=0}^{\infty}x\left(\frac14\right)^x\right)
      \left(\frac34\sum_{y=0}^{\infty}y\left(\frac14\right)^y\right)
      =\frac19.
\end{aligned}
\]
Hence,
\[
\begin{aligned}
    \Exp[Z^2]
    &=\Exp[X^2]+2\Exp[XY]+\Exp[Y^2]
      =\frac43,\\
    \Var(Z)
    &=\Exp[Z^2]-\bigl(\Exp[Z]\bigr)^2
      =\frac43-\left(\frac23\right)^2
      =\frac89.
\end{aligned}
\]
Therefore,
\[
    \boxed{\Exp[Z]=\frac23,\qquad \Var(Z)=\frac89}.
\]
\end{solution}

\begin{exercise}{4.70}
Consider Review Exercise 3.64 on page 127. There are two service lines. The random variables $X$ and $Y$ are the proportions of time that line 1 and line 2 are in use, respectively. The joint probability density function for $(X,Y)$ is given by
\[
    f(x,y)=
    \begin{cases}
        \dfrac32(x^2+y^2), & 0\leq x,y\leq1,\\
        0, & \text{elsewhere}.
    \end{cases}
\]
\begin{exercisesubparts}
    \item Determine whether or not $X$ and $Y$ are independent.
    \item It is of interest to know something about the proportion of $Z=X+Y$, the sum of the two proportions. Find $\Exp[X+Y]$. Also find $\Exp[XY]$.
    \item Find $\Var(X)$, $\Var(Y)$, and $\Cov(X,Y)$.
    \item Find $\Var(X+Y)$.
\end{exercisesubparts}
\end{exercise}

\begin{solution}
\noindent\textup{(a)} The marginal densities are
\[
\begin{aligned}
    f_X(x)
    &=\int_0^1 \frac32(x^2+y^2)\,\dd y
      =\frac{1+3x^2}{2},\\
    f_Y(y)
    &=\int_0^1 \frac32(x^2+y^2)\,\dd x
      =\frac{1+3y^2}{2},
\end{aligned}
\]
for $0\leq x,y\leq1$. Since $f(x,y)\neq f_X(x)f_Y(y)$ in general,
\[
    \boxed{X\text{ and }Y\text{ are not independent}.}
\]

\medskip
\noindent\textup{(b)} By symmetry,
\[
    \Exp[X]=\Exp[Y]
    =\int_0^1 x\frac{1+3x^2}{2}\,\dd x
    =\frac58.
\]
Therefore, by Theorem 4.7,
\[
    \boxed{\Exp[X+Y]=\Exp[X]+\Exp[Y]=\frac54}.
\]
Also,
\[
    \boxed{\Exp[XY]
    =\int_0^1\int_0^1 xy\frac32(x^2+y^2)\,\dd x\,\dd y
    =\frac38}.
\]

\medskip
\noindent\textup{(c)} Again by symmetry,
\[
    \Exp[X^2]=\Exp[Y^2]
    =\int_0^1 x^2\frac{1+3x^2}{2}\,\dd x
    =\frac7{15}.
\]
Therefore,
\[
    \Var(X)=\Var(Y)
    =\frac7{15}-\left(\frac58\right)^2
    =\frac{73}{960},
\]
and, by Theorem 4.4,
\[
    \Cov(X,Y)
    =\frac38-\left(\frac58\right)^2
    =-\frac1{64}.
\]
Thus,
\[
    \boxed{\Var(X)=\Var(Y)=\frac{73}{960},
    \qquad
    \Cov(X,Y)=-\frac1{64}}.
\]

\medskip
\noindent\textup{(d)} Using the moments already found,
\[
\begin{aligned}
    \Exp[(X+Y)^2]
    &=\Exp[X^2]+2\Exp[XY]+\Exp[Y^2]
      =\frac{101}{60},\\
    \Var(X+Y)
    &=\frac{101}{60}-\left(\frac54\right)^2
      =\boxed{\frac{29}{240}}.
\end{aligned}
\]
\end{solution}

\end{document}
